%% kapa: exam=1 round=37 year=2026 period=2 form=A subject=accounting subject_ko=회계학 seq=24 number=64
%% source: Q-Net 감정평가사 기출문제 (cst00309) artlSeq=5258924 file=raw/1008/5258924/2026년 제37회 감정평가사 1차 2교시.pdf
%% conversion: kordoc 4.24.51 -> tools/gen1.py (auto); review=proofread-sonnet; flags=none
\begin{question}{64}
리스에 관한 설명으로 옳은 것을 모두 고른 것은?
\begin{framedtext}
ㄱ. 리스이용자의 규모, 특성, 상황이 서로 다르기 때문에, 기초자산이 소액인지는 상대적 기준에 따라 평가한다. \par
ㄴ. 소액 기초자산 리스에 대한 리스회계처리 선택은 기초자산의 유형별로 적용해야 한다. \par
ㄷ. 리스약정일은 리스계약일과 리스의 주요 조건에 대하여 계약당사자들이 합의한 날 중 이른 날을 말한다. \par
ㄹ. 리스이용자는 리스부채의 원금에 해당하는 현금 지급액은 현금흐름표에 재무활동으로 분류하고, 리스부채 측정치에 포함되지 않은 단기리스료, 소액자산 리스료, 변동리스료는 현금흐름표에 영업활동으로 분류한다. \par
\end{framedtext}
\begin{choices}
  \item ㄱ, ㄴ
  \item ㄱ, ㄹ
  \item ㄴ, ㄷ
  \item ㄴ, ㄹ
  \item ㄷ, ㄹ
\end{choices}
\end{question}
\answer{5}
